\documentclass{article}
\input{../../../../lib.sty}


\title{\texttt{fn cdp\_delta}}
\author{\MichaelShoemateAuthor}

\begin{document}
\maketitle
\contrib

Proves soundness of \texttt{fn cdp\_delta} in \asOfCommit{mod.rs}{0546d1e9a}.
This proof uses the conversion bound from \cite[Proposition~7 and Equation~(2), NeurIPS 2020 version]{CKS20}.

\section{Hoare Triple}
\subsection*{Precondition}
\subsubsection*{Compiler-verified}
\texttt{rho} and \texttt{eps} are of type \texttt{f64}.

\subsubsection*{Caller-verified}
None

\subsection*{Pseudocode}
\lstinputlisting[language=Python,firstline=2]{./pseudocode/cdp_delta.py}

\subsection*{Postcondition}
\begin{theorem}
    For nonnegative privacy parameters $\rho$ and $\epsilon$,
    \texttt{cdp\_delta} either returns an error or returns $\delta$ such that
    any $\rho$-zCDP measurement is also $(\epsilon,\delta)$-differentially private.
\end{theorem}

\begin{proof}
    If $\rho=0$, the measurement is $0$-DP.
    An infinite $\epsilon$ is a vacuous privacy bound.
    Otherwise, if $\rho$ is infinite, $\delta=1$ is sufficient.
    These cases are handled before the search.

    For finite $\rho>0$ and finite $\epsilon\geq 0$, by the cited proposition,
    any finite $\alpha>1$ gives the valid bound
    \[
        \delta(\alpha)=
        \frac{\exp((\alpha-1)(\alpha\rho-\epsilon)+\alpha\ln(1-1/\alpha))}
             {\alpha-1}.
    \]
    The search starts with $\texttt{a\_max}\geq 2$ and $\texttt{a\_min}=1.01$,
    and updates the bounds only to their midpoint.
    Thus, on success, $\alpha=\texttt{a\_max}>1$.
    Optimality of the selected order is not required for privacy.

    The final calculation upper bounds $\alpha\rho-\epsilon$ by directed arithmetic.
    The rounding direction for $\alpha-1$ is chosen according to the sign of this upper bound,
    so \texttt{t1} upper bounds $(\alpha-1)(\alpha\rho-\epsilon)$.
    Rounding $1/\alpha$ down makes $\ln(1-1/\alpha)$ larger,
    and the remaining upward-rounded operations make \texttt{t2} an upper bound on $\alpha\ln(1-1/\alpha)$.
    The numerator is rounded up and the positive denominator down.
    Finally, taking the minimum with 1 preserves the privacy guarantee.
\end{proof}

\bibliographystyle{alpha}
\bibliography{mod}

\end{document}
